% Generado por scripts/70_tablas_figuras.R a partir de resultados-v1. No editar a mano.
\begin{table}[htbp]
\centering
\small
\caption{Tabla 6.13—Construcción de la muestra}
\label{tab:tabla_6_13}
\begin{threeparttable}
\begin{tabular}{@{}lccc@{}}
\toprule
 & (1) & (2) & (3) \\
 & Proyectos o tratados & Controles & Observaciones \\
\midrule
\multicolumn{4}{@{}l}{\textit{Panel A. Proyectos de Verra}} \\
Proyectos de Verra en Colombia, todos los sectores (Tabla 1) & 101 &  &  \\
Con actividad AFOLU REDD+ declarada & 45 &  &  \\
Menos: punto en la sede del proponente & 1 &  &  \\
Con ubicación válida (Figura 3 y escalera de estado) & 44 &  &  \\
De ellos, con año de inicio publicado & 44 &  &  \\
De ellos, con polígono en Karnik et al. (2025) & 19 &  &  \\
De ellos, sin tratamiento (metodología o área) & 2 &  &  \\
Fuera de los 45: metodología REDD+ sin la actividad & 1 &  &  \\
\addlinespace
\multicolumn{4}{@{}l}{\textit{Panel B. Universo de la Tabla A.3 (registros)}} \\
De Verra: los 44 válidos y el de metodología sin la actividad & 45 &  &  \\
Polígonos de deforestación evitada de otros registros & 14 &  &  \\
Total de registros & 59 &  &  \\
Proyectos únicos & 58 &  &  \\
\addlinespace
\multicolumn{4}{@{}l}{\textit{Panel C. Base de polígonos (Karnik et al., 2025)}} \\
Polígonos en Colombia (todos los registros y tipos) & 55 &  &  \\
De deforestación evitada & 32 &  &  \\
De ellos, de proyectos de Verra & 18 &  &  \\
Municipios que cubren ($\geq$ 1 \% y $\geq$ 100 ha) & 68 &  &  \\
\addlinespace
\multicolumn{4}{@{}l}{\textit{Panel D. Municipios}} \\
Municipios del panel & 1.122 &  & 26.928 \\
Tratados (definición principal) & 42 &  &  \\
De ellos, con adopción fuera de 2013--2019 & 6 &  &  \\
Excluidos: único proyecto retirado o rechazado & 6 &  &  \\
Nunca tratados & 1.074 &  &  \\
Menos: con vínculo con proyectos de carbono & 88 &  &  \\
Controles limpios & 986 &  &  \\
Tratados de la muestra principal & 36 &  &  \\
De ellos, con proyecto REDD+ de Verra & 25 &  &  \\
De ellos, solo con proyectos de RENARE o Cercarbono & 11 &  &  \\
De ellos, con ubicación verificada & 26 &  &  \\
De ellos, con su polígono en otros municipios & 2 &  &  \\
De ellos, sin polígono en la base & 8 &  &  \\
\addlinespace
\multicolumn{4}{@{}l}{\textit{Panel E. Muestra de cada especificación}} \\
Binaria principal (Tabla 4) & 36 & 986 & 24.528 \\
Ubicación verificada (Tabla 4, columna 6) & 26 & 986 & 24.288 \\
Tratamiento continuo y gradiente (Tabla 5) & 26 & 986 & 24.288 \\
Dosis híbrida (Tabla 5) & 30 & 986 & 24.384 \\
Reasignación geográfica (Tabla 10) & 72 & 953 & 24.600 \\
Adopción según el polígono (Tabla 10) & 34 & 986 & 24.480 \\
Evaluación por polígonos (Tabla 9), proyectos & 25 &  &  \\
Regresión discontinua (Tabla A.4), proyectos & 26 &  &  \\
\bottomrule
\end{tabular}
\begin{tablenotes}[flushleft]
\footnotesize
\item \textit{Notas:} Todas las cifras salen de las bases y de las salidas finales (script 78). Unidad: proyecto único por identificador en los paneles A y C, registro en el B y municipio en el D y el E, con observaciones municipio-año del panel balanceado 2001--2024; la evaluación por polígonos y la regresión discontinua cuentan proyectos y usan celdas de 2 × 2 km. Panel A: REDD+ es la actividad AFOLU declarada en el registro, y el tratamiento exige además una metodología REDD+. El excluido es CO2ROZO, cuyo punto corresponde a la sede del proponente. Los dos válidos sin tratamiento son el de corredores de jaguar, cuya metodología es de restauración, y el corredor de robles, cuya área declarada no puede asignarse al municipio de su punto (capítulo 4). Yaguara Llanos del Yarí no está entre los 45: usa una metodología REDD+ sin la actividad declarada, fue retirado y no tiene fecha de inicio, y su municipio, Calamar, sale de la muestra por tener solo proyectos retirados. Panel B: los 59 registros corresponden a 58 proyectos porque Chonta Corozo figura en Verra y en EcoRegistry. Panel D: los tratados combinan Verra, RENARE y Cercarbono; la adopción es el año de inicio más temprano entre sus proyectos REDD+ no retirados ni rechazados, y la ubicación se verifica con los polígonos de deforestación evitada. Los conteos de Cercarbono, en registros y proyectos, y del RENARE, en iniciativas, están en el capítulo 4.
\item \textit{Fuente:} Verra; Karnik et al. (2025); EcoRegistry; BioCarbon Registry; RENARE; Cercarbono; cálculos propios.
\end{tablenotes}
\end{threeparttable}
\end{table}
